% ============================================================================
%  matrix_mechanics.tex
%  VSEPR-SIM / thenewmethods
%
%  Document 2 of 3 -- The New Methods
%  Matrix mechanics: forces, moments, torsion, stiffness,
%  least-squares, and the Z-bead extension
% ============================================================================

\documentclass[11pt, letterpaper]{article}

\usepackage{amsmath, amssymb, bm}
\usepackage{geometry}
\usepackage{booktabs}
\usepackage{array}
\usepackage{xcolor}
\usepackage{microtype}
\usepackage{parskip}
\usepackage{hyperref}

\geometry{margin=1.1in}

\definecolor{accent}{RGB}{30, 80, 160}

\hypersetup{colorlinks=true, linkcolor=accent, urlcolor=accent}

\title{\textbf{The New Methods II}\\[0.4em]
\large Matrix Mechanics: Forces, Moments, Torsion,\\
Stiffness, Least Squares, and the $Z$-Bead Extension}
\author{VSEPR-SIM thenewmethods}
\date{}

\usepackage{titlesec}
\titleformat{\section}{\large\bfseries\color{accent}}{}{0em}{}[\vspace{-0.3em}\rule{\linewidth}{0.6pt}]
\titleformat{\subsection}{\normalsize\bfseries}{}{0em}{}

% ── macros ──────────────────────────────────────────────────────────────────
\newcommand{\uvec}[1]{\hat{\mathbf{#1}}}
\newcommand{\bL}{\mathbf{L}}
\newcommand{\bA}{\mathbf{A}}
\newcommand{\bP}{\mathbf{P}}
\newcommand{\bK}{\mathbf{K}}
\newcommand{\bu}{\mathbf{u}}
\newcommand{\bx}{\mathbf{x}}
\newcommand{\bv}{\mathbf{v}}
\newcommand{\bF}{\mathbf{F}}
\newcommand{\bZ}{\mathbf{Z}}
\newcommand{\bq}{\mathbf{q}}
\newcommand{\argmin}{\operatorname{arg\,min}}

% ============================================================================
\begin{document}
\maketitle
\thispagestyle{empty}

\noindent
Everything below is undergraduate mechanics in different costumes, unified
by a single observation: write the unknowns as a vector, build a system
$\bA\bx = \mathbf{b}$, and solve.

% ────────────────────────────────────────────────────────────────────────────
\section{Three-Force System}
% ────────────────────────────────────────────────────────────────────────────

Unknown magnitudes $L_6, L_{12}, L_{18}$ along known unit directions
$\hat{u}_1, \hat{u}_2, \hat{u}_3$.
Equilibrium against external load $\mathbf{P}$:

\[
\underbrace{\begin{bmatrix}
u_{1x} & u_{2x} & u_{3x}\\
u_{1y} & u_{2y} & u_{3y}\\
u_{1z} & u_{2z} & u_{3z}
\end{bmatrix}}_{\bA}
\underbrace{\begin{bmatrix} L_6 \\ L_{12} \\ L_{18} \end{bmatrix}}_{\bL}
=
\underbrace{\begin{bmatrix} P_x \\ P_y \\ P_z \end{bmatrix}}_{\bP}
\]

\medskip\noindent
\textbf{Ideal (exact geometry):}
\[
\bL = \bA^{-1}\bP
\]

\noindent
\textbf{Real (noisy / overdetermined):}
\[
\bL_{\mathrm{real}} = (\bA^T\bA)^{-1}\bA^T\bP
\]

\noindent
This covers cable tensions, strut forces, rod loads, ideal joint load sharing.

% ────────────────────────────────────────────────────────────────────────────
\section{Add Moments About a Point}
% ────────────────────────────────────────────────────────────────────────────

Each force $L_k\hat{u}_k$ has position vector $\mathbf{r}_k$ from point $O$.
Stack force and moment equations:

\[
\begin{bmatrix}
  \hat{u}_1 & \hat{u}_2 & \hat{u}_3 \\[4pt]
  \mathbf{r}_1\times\hat{u}_1 & \mathbf{r}_2\times\hat{u}_2 & \mathbf{r}_3\times\hat{u}_3
\end{bmatrix}
\begin{bmatrix} L_6 \\ L_{12} \\ L_{18} \end{bmatrix}
=
\begin{bmatrix} \bP \\ \mathbf{M}_{\mathrm{req}} \end{bmatrix}
\]

\noindent
(Each row of the upper block is a $3$-vector; the full matrix is $6\times3$.)
Solve with the pseudo-inverse for a wrench balance.

% ────────────────────────────────────────────────────────────────────────────
\section{Coordinate-Aligned Ideal vs.\ Real Example}
% ────────────────────────────────────────────────────────────────────────────

Assume members aligned with coordinate axes:
$\hat{u}_1=(1,0,0)^T$, $\hat{u}_2=(0,1,0)^T$, $\hat{u}_3=(0,0,1)^T$, and load
$\bP=(120,\;80,\;-40)^T$ N. Then $\bA = \mathbf{I}$, so

\[
\bL_{\mathrm{ideal}} = (120,\;80,\;-40)^T \text{ N}
\]

\noindent
Negative sign: member $18$ is in compression.

\medskip\noindent
\textbf{Real case} -- directions slightly off:

\[
\bA = \begin{bmatrix}
0.98 & 0.05 & 0.00\\
0.04 & 0.99 & 0.10\\
0.00 & 0.03 & 0.95
\end{bmatrix}
\]

Solve $\bA\bL=\bP$ for adjusted internal loads.
Geometry error changes load sharing: what reality loves to do precisely
when your design margin is already garbage.

% ────────────────────────────────────────────────────────────────────────────
\section{Stiffness Method}
% ────────────────────────────────────────────────────────────────────────────

Connect force and displacement directly:
\[
\bK\bu = \mathbf{f}
\]

\noindent
$\bK$ = stiffness matrix, $\bu$ = nodal displacement vector,
$\mathbf{f}$ = nodal force vector.

For a 1D two-spring system:
\[
\begin{bmatrix}
k_1+k_2 & -k_2 \\
-k_2    & k_2+k_3
\end{bmatrix}
\begin{bmatrix} u_1 \\ u_2 \end{bmatrix}
=
\begin{bmatrix} F_1 \\ F_2 \end{bmatrix}
\]

Gateway to FEA; yields beam deflection, truss extension, support reactions.

% ────────────────────────────────────────────────────────────────────────────
\section{Least-Squares / Pseudo-Inverse}
% ────────────────────────────────────────────────────────────────────────────

Overdetermined or noisy system:
\[
\bA\bx \approx \mathbf{b}
\qquad\Longrightarrow\qquad
\bx = \bA^+\mathbf{b}
\]

Residual and residual norm:
\[
\mathbf{r} = \mathbf{b} - \bA\bx \qquad \|\mathbf{r}\|
\]

Useful for force reconstruction from sensor data, best-fit reaction estimates,
and real-vs-ideal comparison.

% ────────────────────────────────────────────────────────────────────────────
\section{Torsional System}
% ────────────────────────────────────────────────────────────────────────────

Torsion identity:
\[
\frac{T}{J} = \frac{\tau}{r} = \frac{G\theta}{L}
\]

Torsional stiffness of segment $i$:
\[
k_i^{(t)} = \frac{G_i J_i}{L_i}
\]

Element stiffness matrix for a two-node torsional element:
\[
k_i^{(t)}\begin{bmatrix}1 & -1\\-1 & 1\end{bmatrix}
\]

Assembled shaft system:
\[
\bK_\theta\,\boldsymbol{\theta} = \mathbf{T}
\]

Post-solve, shear stress and angle of twist in each segment:
\[
\tau_{\max,i} = \frac{T_i c_i}{J_i}
\qquad\qquad
\phi_i = \frac{T_i L_i}{G_i J_i}
\]

% ────────────────────────────────────────────────────────────────────────────
\section{Unified Block System}
% ────────────────────────────────────────────────────────────────────────────

Define the combined unknown vector:
\[
\bq =
\begin{bmatrix} L_6 \\ L_{12} \\ L_{18} \\ \theta_1 \\ \theta_2 \\ \theta_3 \end{bmatrix}
\]

One block equation:
\[
\begin{bmatrix}
\bA_f & \mathbf{0} \\
\bA_m & \mathbf{0} \\
\mathbf{0} & \bK_\theta
\end{bmatrix}
\bq
=
\begin{bmatrix} \bP \\ \mathbf{M} \\ \mathbf{T} \end{bmatrix}
\]

$\bA_f$ = force balance, $\bA_m$ = moment balance, $\bK_\theta$ = torsion.
One framework for force members, moment equilibrium, and rotational response.

% ────────────────────────────────────────────────────────────────────────────
\section{Compact $L_{6,12,18}$ Notation Summary}
% ────────────────────────────────────────────────────────────────────────────

\[
\bL = \begin{bmatrix} L_6 \\ L_{12} \\ L_{18} \end{bmatrix}
\qquad
\bA\bL = \bP
\]

\[
\bL_{\mathrm{ideal}} = \bA^{-1}\bP
\qquad
\bL_{\mathrm{real}} = (\bA^T\bA)^{-1}\bA^T\bP
\qquad
\mathbf{r} = \bP - \bA\bL_{\mathrm{real}}
\]

Exact solution, best-fit solution, error measure. Neat.

% ────────────────────────────────────────────────────────────────────────────
\section{$Z$-Bead Extension: Forces from First Principles}
% ────────────────────────────────────────────────────────────────────────────

Each particle $i$ carries state:
\[
Z_i =
\begin{bmatrix} \bx_i \\ \bv_i \\ s_i \\ q_i \\ m_i \\ \varepsilon_i \end{bmatrix}
\qquad
\mathbf{Z} = \{Z_1, Z_2, \ldots, Z_N\}
\]

Force system becomes $Z$-dependent:
\[
\bA(\mathbf{Z})\,\bL = \bP(\mathbf{Z})
\]

Pairwise interaction:
\[
\bF_{ij} = f(s_i, s_j, r_{ij}, q_i, q_j, \nabla\Phi)
\qquad r_{ij} = \|\bx_i - \bx_j\|
\]

Total force on bead $i$:
\[
\bF_i = \sum_{j\neq i} \bF_{ij}
\]

\subsection*{Macro Property Extraction}

\[
\rho = \frac{\sum_i m_i}{V(\mathbf{Z})}
\qquad
\boldsymbol{\sigma} = \frac{1}{V}\sum_i \bF_i \otimes \bx_i
\qquad
E = \sum_i \varepsilon_i + \sum_{i<j} U_{ij}
\]

\subsection*{Ideal vs.\ Real in $Z$-Language}

\[
\text{Ideal: } \mathbf{Z}_{\mathrm{ideal}} \;\Rightarrow\; \bA(\mathbf{Z})^{-1}\bP(\mathbf{Z})
\]
\[
\text{Real: } \bL_{\mathrm{real}} = (\bA^T\bA)^{-1}\bA^T\bP(\mathbf{Z}_{\mathrm{real}})
\qquad
\mathbf{r} = \bP - \bA(\mathbf{Z})\bL
\]

\subsection*{Coarse-Grained Force Channels}

Each channel $L_k(\mathbf{Z})$ is a projection of many bead forces:
\[
L_k(\mathbf{Z}) = \sum_{i\in\Omega_k} \bF_i \cdot \hat{u}_k
\]

\subsection*{Unified System}

\[
\begin{cases}
\bF(\mathbf{Z}) = \bA(\mathbf{Z})\bL \\[4pt]
\mathbf{T}(\mathbf{Z}) = \bK_\theta(\mathbf{Z})\,\boldsymbol{\theta} \\[4pt]
\mathbf{P}(\mathbf{Z}) \to \text{macro properties}
\end{cases}
\]

% ────────────────────────────────────────────────────────────────────────────
\section{Advanced Methods from the Same Framework}
% ────────────────────────────────────────────────────────────────────────────

\begin{tabular}{@{}ll@{}}
\toprule
Goal & Method \\
\midrule
Natural modes / stability & $\bK\mathbf{v} = \lambda\mathbf{v}$ (eigenvalue problem) \\
Best-fit load estimate & $\min_{\bx}\|\bA\bx - \mathbf{b}\|^2$ \\
Design under equilibrium & $\min J(\bx)$ subject to $\bA\bx = \mathbf{b}$ \\
Numerical stability check & $\kappa(\bA)$ (condition number) \\
Sensor fusion & $\bL = \bA^+\bP_{\mathrm{sensor}}$ \\
\bottomrule
\end{tabular}

\bigskip\noindent
From force balance through eigenvalue stability through bead dynamics,
it is the same algebraic skeleton wearing different physics.

\end{document}
